% ==========================================================================
%  Chapter 3 — Technical Foundations
%  User-Mode Linux: A Comprehensive Technical Report
% ==========================================================================

\chapter{Technical Foundations}
\label{ch:foundations}

\epigraph{Any sufficiently complex virtual machine architecture is
indistinguishable from a real one.}{---with apologies to Arthur C. Clarke}

This chapter establishes the technical concepts required to understand
how a kernel can run as an ordinary userspace process.  We begin with
the classical virtualization taxonomy and show where user-mode kernels
fit---and where they do not.  We then trace the x86 privilege model
from its original ring-based design through the hardware virtualization
extensions that now underpin \KVM{}.  Finally, we examine the Linux
process-model primitives that \UML{} exploits---\syscall{ptrace},
\Seccomp{}, \syscall{mmap}, and the \KVM{} ioctl interface---and compare
the three trap mechanisms that have defined \UML{}'s performance
envelope over the past quarter century.


% ------------------------------------------------------------------
\section{Virtualization Taxonomy}
\label{sec:foundations:taxonomy}

\subsection{The Popek--Goldberg Requirements}

The formal study of virtual machines begins with Goldberg's 1973
taxonomy of \emph{virtual machine monitors} (VMMs)~\cite{goldberg1973}
and the celebrated Popek--Goldberg theorem of 1974~\cite{popek1974}.
Their framework identifies three properties a VMM must satisfy:

\begin{enumerate}[label=\textbf{P\arabic*.},leftmargin=2.5em]
  \item \textbf{Equivalence (Fidelity).}
        A program running under the VMM exhibits behavior identical to
        that observed when running directly on the hardware, with at
        most timing differences.
  \item \textbf{Resource Control (Safety).}
        The VMM retains complete control of the virtualized resources.
        No guest operation may access a resource not explicitly
        allocated to that guest.
  \item \textbf{Efficiency (Performance).}
        A statistically dominant fraction of guest instructions execute
        directly on the hardware without VMM intervention.
\end{enumerate}

The theorem proves that a VMM satisfying all three properties can be
constructed for any \emph{third-generation architecture}\footnote{
  Popek and Goldberg use ``third generation'' to mean a conventional
  von~Neumann machine with a supervisor/user mode split, exemplified by
  IBM~System/370.}
if and only if the set of \emph{sensitive instructions}\footnote{
  Instructions whose behavior depends on the current privilege level or
  that affect the machine's resource-allocation state.}
is a subset of the set of \emph{privileged instructions} (those that
trap when executed in user mode).  When every sensitive instruction is
privileged, the VMM can intercept all guest attempts to observe or
modify machine state by running the guest at a lower privilege level
and catching the resulting traps.

\subsection{Taxonomy of Hypervisors}

Smith and Nair~\cite{smith2005book} refine the taxonomy into two
orthogonal dimensions:

\begin{description}[style=nextline]
  \item[Type~1 (bare-metal) vs.\ Type~2 (hosted) hypervisors.]
    A Type~1 hypervisor runs directly on the hardware with no
    intervening operating system (e.g., VMware ESXi, Xen).  A Type~2
    hypervisor runs as a process atop a host OS (e.g., VMware
    Workstation, VirtualBox).  \KVM{} blurs the boundary: it is a
    loadable module that turns the Linux kernel itself into a Type~1
    hypervisor while simultaneously serving as a host OS.

  \item[Full virtualization vs.\ paravirtualization vs.\ OS-level.]
    Full virtualization presents an unmodified hardware interface to
    the guest; the guest kernel need not know it is virtualized.
    Paravirtualization (exemplified by Xen~\cite{barham2003}) modifies
    the guest kernel to replace sensitive instructions with explicit
    \emph{hypercalls}.  OS-level virtualization (containers) shares a
    single kernel instance and isolates guests via namespaces and
    cgroups~\cite{soltesz2007}.
\end{description}

\subsection{Where User-Mode Kernels Fit}

\UML{} does not fit cleanly into any of these categories, and this
taxonomic awkwardness is part of what makes it interesting.

\begin{itemize}
  \item It is not a hypervisor in the Popek--Goldberg sense, because
        it does not satisfy the Efficiency property: in the original
        \syscall{ptrace}-based design, \emph{every} guest system call
        requires VMM intervention.
  \item It is not OS-level virtualization, because it runs an
        independent kernel instance with its own scheduler, memory
        management, and driver stack.
  \item It is not paravirtualization, because the guest userspace is
        unmodified---ordinary Linux binaries run without recompilation.
  \item With the \KVM{} backend (Chapter~\ref{ch:kvm}), \UML{}
        gains the Efficiency property: a dominant fraction of guest
        instructions execute directly on the CPU inside a \KVM{}
        virtual machine, and the LSTAR gadget handles trivial system
        calls entirely in-guest without any VMM exit.
\end{itemize}

Perhaps the most accurate characterization is that \UML{} is a
\emph{process-level system virtual machine}~\cite{smith2005book}:
it implements a complete system interface (an entire Linux kernel)
within the confines of a host process.  The \KVM{} backend promotes
it to a \emph{hybrid} system VM that uses hardware virtualization for
the hot path while retaining the hosted-process model for I/O,
lifecycle management, and debugging.


% ------------------------------------------------------------------
\section{The x86 Privilege Model}
\label{sec:foundations:privilege}

\subsection{Rings, CPL, and DPL}

The x86 architecture defines four privilege levels, conventionally
called \emph{rings} 0 through 3.  The \textbf{Current Privilege Level}
(CPL) is the ring at which the processor is currently executing; it is
encoded in the two least-significant bits of the CS and SS segment
selectors.  The \textbf{Descriptor Privilege Level} (DPL) field in
segment descriptors, gate descriptors, and page-table entries gates
access: a code segment can only be called from a CPL that is at least
as privileged as the target DPL.

\begin{figure}[ht]
\centering
\begin{tikzpicture}[
  ring/.style={draw, thick, circle, minimum size=#1, color=UMLBlue!80},
  rlabel/.style={font=\small\bfseries, color=UMLInk},
  desc/.style={font=\footnotesize, color=UMLMuted, text width=3.5cm, align=center}
]
  % Rings (outermost first for layering)
  \node[ring=7.0cm, fill=UMLBlue!5]  (r3) at (0,0) {};
  \node[ring=5.2cm, fill=UMLTeal!5]  (r2) at (0,0) {};
  \node[ring=3.4cm, fill=UMLGreen!5] (r1) at (0,0) {};
  \node[ring=1.6cm, fill=KernelGold!15] (r0) at (0,0) {};

  % Labels
  \node[rlabel] at (0,0)    {Ring 0};
  \node[rlabel] at (0,1.1)  {Ring 1};
  \node[rlabel] at (0,2.0)  {Ring 2};
  \node[rlabel] at (0,2.9)  {Ring 3};

  % Descriptions
  \node[desc] at (4.2, 0)    {Kernel mode\\(traditional OS)};
  \node[desc] at (4.2, 1.1)  {Device drivers\\(rarely used)};
  \node[desc] at (4.2, 2.0)  {Device drivers\\(rarely used)};
  \node[desc] at (4.2, 2.9)  {User mode\\(\textbf{UML runs here})};

  % Brace for UML
  \draw[decorate, decoration={brace, amplitude=6pt, raise=4pt},
        thick, UMLRed]
    (5.8, 2.3) -- (5.8, 3.5)
    node[midway, right=12pt, font=\footnotesize\bfseries, color=UMLRed,
         text width=2.5cm, align=left]
    {UML kernel +\\UML userspace};
\end{tikzpicture}
\caption{The x86 privilege ring model.  Traditional kernels occupy
Ring~0; \UML{} runs entirely at Ring~3 (CPL=3), sharing the same
privilege level as ordinary user processes.}
\label{fig:x86-rings}
\end{figure}

Traditional operating systems use a two-ring model: the kernel at
Ring~0, user processes at Ring~3.  Rings~1 and~2 exist in hardware
but are almost universally unused.\footnote{OS/2 and a few
microkernel-inspired designs briefly used Ring~1 or Ring~2 for device
drivers, but the x86-64 long mode effectively reduces the model to
two rings in practice.}

\subsection{Implications for UML}

\UML{} runs entirely at Ring~3.  This has several consequences:

\begin{enumerate}
  \item \textbf{No privileged instructions.}
        The \UML{} kernel cannot execute \texttt{cli}/\texttt{sti}
        (interrupt control), \texttt{mov~cr3} (page-table switching),
        \texttt{lgdt}/\texttt{lidt} (descriptor-table loading), or
        any other privileged instruction.  Every such operation must
        be emulated via host system calls.

  \item \textbf{No direct hardware access.}
        All I/O must pass through the host kernel's system-call
        interface.  \UML{}'s block driver issues \syscall{read}/\syscall{write}
        on a host file descriptor; its network driver uses host sockets
        or TUN/TAP devices.

  \item \textbf{Shared address space with the host.}
        The \UML{} process occupies the same 48-bit (or 57-bit) virtual
        address space managed by the host kernel.  The host kernel's
        own mappings (above the canonical hole) are inaccessible from
        CPL=3, but the host's ASLR layout, vDSO, and stack placement
        all constrain where \UML{} can place its own data structures.

  \item \textbf{System call interception required.}
        When guest userspace issues a \texttt{syscall} instruction, it
        would normally trap into the \emph{host} kernel, not the
        \UML{} kernel.  \UML{} must intercept this trap and redirect
        it to its own system-call handler.  The mechanism for this
        interception---\syscall{ptrace}, \Seccomp{}, or \KVM{}---is
        the defining architectural choice of each \UML{} backend.
\end{enumerate}


% ------------------------------------------------------------------
\section{Hardware Virtualization Extensions}
\label{sec:foundations:hwvirt}

The x86 architecture famously violates the Popek--Goldberg
requirements: several sensitive instructions (notably \texttt{popf},
\texttt{sgdt}, and \texttt{pushf}) are not privileged and therefore
do not trap when executed in user mode~\cite{adams2006}.  This made
software-only full virtualization on x86 difficult, requiring either
binary translation (VMware) or paravirtualization (Xen).

Hardware virtualization extensions---Intel VT-x (2005) and AMD-V
(2006)---solve this by adding a new, orthogonal dimension of privilege.

\subsection{Intel VT-x (VMX)}

VT-x~\cite{intel-sdm} introduces \textbf{VMX root mode} (the
hypervisor) and \textbf{VMX non-root mode} (the guest).  The
transition between the two is mediated by a per-vCPU data structure
called the \textbf{Virtual Machine Control Structure} (VMCS):

\begin{description}[style=nextline]
  \item[VM Entry (\texttt{VMLAUNCH} / \texttt{VMRESUME}).]
    The hypervisor loads guest state from the VMCS and transfers
    control to the guest.  The guest runs at whatever CPL the VMCS
    guest-state area specifies---typically Ring~0 for a full-system
    guest, or Ring~3 for a process-level guest like \UML{}'s \KVM{}
    backend.

  \item[VM Exit.]
    Certain events cause the processor to stop guest execution,
    save guest state into the VMCS, and transfer control back to
    VMX root mode.  The exit reason is recorded in the VMCS.
    Events that cause exits include: execution of privileged
    instructions, I/O instructions (if the I/O bitmap so indicates),
    external interrupts, page faults (if configured), and the
    \texttt{CPUID} instruction.

  \item[Extended Page Tables (EPT).]
    EPT provides a second level of address translation: guest-physical
    addresses (GPAs) are translated to host-physical addresses (HPAs)
    by hardware, without any hypervisor intervention for non-faulting
    accesses.  This eliminates the ``shadow page table'' overhead that
    made pre-EPT virtualization expensive.
\end{description}

\subsection{AMD-V (SVM)}

AMD's Secure Virtual Machine~\cite{amd-apm} follows the same
principles with different terminology: the per-vCPU structure is the
\textbf{Virtual Machine Control Block} (VMCB); the hardware nested
paging feature is called \textbf{Nested Page Tables} (NPT); and the
entry/exit instructions are \texttt{VMRUN} and---implicitly---any
intercepted event.  The functional model is equivalent for the
purposes of this report.

\subsection{KVM and the ioctl Interface}

\KVM{}~\cite{kivity2007} exposes hardware virtualization to userspace
through a layered ioctl interface on \texttt{/dev/kvm}:

\begin{enumerate}
  \item \texttt{KVM\_CREATE\_VM} returns a VM file descriptor.
  \item \texttt{KVM\_SET\_USER\_MEMORY\_REGION} on the VM fd maps
        host-virtual memory into the guest-physical address space
        (backing EPT/NPT).
  \item \texttt{KVM\_CREATE\_VCPU} on the VM fd returns a vCPU file
        descriptor, with a shared \texttt{struct kvm\_run} page
        accessible via \syscall{mmap} on the vCPU fd.
  \item \texttt{KVM\_RUN} on the vCPU fd enters VMX/SVM non-root mode.
        When the guest causes an exit, \texttt{KVM\_RUN} returns and
        the exit reason + auxiliary data are available in the shared
        \texttt{kvm\_run} page.
\end{enumerate}

\subsection{Why This Matters for UML}

\UML{}'s \KVM{} backend uses this interface to run guest code at
Ring~3 \emph{inside a virtual machine}, while the \UML{} kernel
process remains an ordinary userspace process on the host.  The
critical insight is that \KVM{}'s VM entry/exit cycle is far cheaper
than a \syscall{ptrace} round-trip or even a \Seccomp{} SIGSYS
delivery.  A \texttt{KVM\_RUN} $\rightarrow$ \texttt{KVM\_EXIT\_IO}
round-trip costs approximately \SIrange{100}{150}{\nano\second}; with
the LSTAR gadget (Section~\ref{sec:foundations:trap-compare}), trivial
system calls complete in approximately 90 cycles without exiting the
VM at all.


% ------------------------------------------------------------------
\section{The Linux Process Model}
\label{sec:foundations:process-model}

From the host kernel's perspective, a running \UML{} instance is an
ordinary multi-threaded process (or set of processes, depending on the
backend).  \UML{} exploits a specific set of Linux system calls to
implement its virtualization layer.

\subsection{ptrace: The Original Trap Mechanism}

The \syscall{ptrace} system call~\cite{ptrace-man} allows one process
(the \emph{tracer}) to observe and control the execution of another
(the \emph{tracee}).  \UML{}'s original backend uses:

\begin{description}[style=nextline]
  \item[\texttt{PTRACE\_SYSEMU}.]
    Resumes the tracee but intercepts the next system call
    \emph{without executing it}.  The tracee stops with a
    \texttt{SIGTRAP|0x80} signal.  This is the key extension that
    Jeff Dike contributed to the Linux kernel specifically for
    \UML{}: without it, the host kernel would execute the guest's
    system call before \UML{} could intercept it.

  \item[\texttt{PTRACE\_GETREGS} / \texttt{PTRACE\_SETREGS}.]
    Read and write the tracee's general-purpose registers.  \UML{}
    reads the registers to determine which system call the guest
    requested (from \texttt{RAX}) and its arguments (from
    \texttt{RDI}, \texttt{RSI}, \texttt{RDX}, \texttt{R10},
    \texttt{R8}, \texttt{R9} per the x86-64 ABI), then writes the
    return value back into \texttt{RAX}.
\end{description}

\subsection{seccomp: The Modern Trap Mechanism}

The \Seccomp{} BPF filter~\cite{seccomp-bpf} was originally designed
for sandboxing (Chrome, Firefox, container runtimes), but its
\texttt{SECCOMP\_RET\_TRAP} return action proved ideal for \UML{}.
When a seccomp filter returns \texttt{SECCOMP\_RET\_TRAP}, the kernel
delivers a \texttt{SIGSYS} signal to the thread that issued the
system call---without executing it.  This is functionally equivalent
to \texttt{PTRACE\_SYSEMU} but with dramatically lower overhead
because the signal is delivered in-process rather than requiring a
context switch to an external tracer.

\subsection{futex: Stub Synchronization}

The \syscall{futex} system call provides efficient userspace-to-kernel
synchronization~\cite{futex-lwn}.  In the \Seccomp{} backend, the
stub process and the \UML{} kernel use a shared \texttt{futex} word
(in the \texttt{stub\_data} page) to implement a lightweight turnstile:
the stub writes its trap state, issues \texttt{FUTEX\_WAKE}, and blocks
on \texttt{FUTEX\_WAIT}; the \UML{} kernel wakes, processes the trap,
writes the result, and wakes the stub.

\subsection{mmap / munmap / mprotect: Guest Memory}

\UML{} implements guest ``physical'' memory as a host-side anonymous
mapping or a file-backed mapping (typically on \texttt{tmpfs}).  Guest
page-table changes are realized by calling \syscall{mmap},
\syscall{munmap}, and \syscall{mprotect} on the host to adjust the
stub process's address space.  Each guest TLB flush translates to a
batch of these host system calls.

\subsection{clone / fork: Guest Processes}

Each guest process in \UML{} is backed by a host process (or thread,
depending on the configuration).  The \syscall{clone} system call
creates the host-side entity; \texttt{CLONE\_VM} may or may not be
used depending on whether the backend shares address spaces.

\subsection{KVM ioctl Interface}
\label{sec:foundations:kvm-api}

The \KVM{} backend~\cite{kvm-api} replaces the \syscall{ptrace} and
\Seccomp{} trap mechanisms with hardware-assisted virtualization.  The
key ioctls are:

\begin{lstlisting}[language=KernelC,caption={KVM ioctl sequence for
UML's v2 backend.},label=lst:kvm-ioctls]
/* VM creation (once per UML instance) */
vm_fd = ioctl(kvm_fd, KVM_CREATE_VM, 0);

/* Map guest physical memory (host tmpfs -> guest GPA) */
struct kvm_userspace_memory_region mem = {
    .slot            = 0,
    .guest_phys_addr = 0,
    .memory_size     = physmem_size,
    .userspace_addr  = (u64)uml_physmem,
};
ioctl(vm_fd, KVM_SET_USER_MEMORY_REGION, &mem);

/* Per-vCPU creation (one per host CPU in the pool) */
vcpu_fd = ioctl(vm_fd, KVM_CREATE_VCPU, cpu_id);
run = mmap(NULL, sizeof(struct kvm_run), PROT_RW,
           MAP_SHARED, vcpu_fd, 0);

/* Hot loop: enter guest, handle exit, re-enter */
while (1) {
    ioctl(vcpu_fd, KVM_RUN, 0);
    switch (run->exit_reason) {
    case KVM_EXIT_IO:
        handle_io_trap(regs, run, vcpu);
        break;
    /* ... other exit reasons ... */
    }
}
\end{lstlisting}


% ------------------------------------------------------------------
\section{Address Space Layout}
\label{sec:foundations:address-space}

\UML{} must carve a complete kernel and user address space out of the
host process's available virtual address range.  The layout differs
between the \Seccomp{} and \KVM{} backends.

\begin{figure}[ht]
\centering
\begin{tikzpicture}[
  block/.style={draw, thick, minimum width=10cm, minimum height=0.8cm,
                font=\footnotesize\ttfamily, anchor=south west},
  label/.style={font=\footnotesize, color=UMLMuted, anchor=east},
  >=Stealth
]
  % Seccomp layout
  \node[font=\small\bfseries, color=UMLBlue] at (5, 7.8) {Seccomp Backend Address Space};

  \node[block, fill=UMLRed!10]    (hk) at (0, 7)
    {Host kernel (inaccessible from CPL=3)};
  \node[label] at (-0.2, 7.4) {\texttt{0xffff...}};

  \node[block, fill=UMLMuted!10]  (gap1) at (0, 6)
    {(canonical hole --- unmapped)};

  \node[block, fill=UMLBlue!10]   (ukt) at (0, 5)
    {UML kernel text + data + bss};
  \node[label] at (-0.2, 5.4) {\texttt{0x6000...}};

  \node[block, fill=UMLGreen!10]  (phys) at (0, 4)
    {UML ``physical memory'' (mmap'd tmpfs / anon)};
  \node[label] at (-0.2, 4.4) {\texttt{uml\_physmem}};

  \node[block, fill=KernelGold!10] (guest) at (0, 3)
    {Guest userspace mappings (via stub mmap)};
  \node[label] at (-0.2, 3.4) {\texttt{0x0000...}};

  \node[block, fill=UMLAmber!10]  (stub) at (0, 2)
    {Stub pages (code + 2 data pages)};
  \node[label] at (-0.2, 2.4) {\texttt{STUB\_START}};

  \draw[decorate, decoration={brace, amplitude=6pt, raise=4pt},
        thick, UMLTeal]
    (10.2, 2) -- (10.2, 5.8)
    node[midway, right=12pt, font=\footnotesize, color=UMLTeal,
         text width=2.5cm, align=left]
    {Managed by\\UML kernel};

  % KVM layout annotation
  \node[font=\small\bfseries, color=UMLBlue] at (5, 1.0) {KVM Backend: adds PML4 guest page tables};
  \node[font=\footnotesize, color=UMLMuted, text width=10cm, align=center] at (5, 0.3)
    {Guest runs in KVM with its own PML4. PML4[448] maps the kernel-half\\
     trampoline (LSTAR, IDT, GDT, TSS, IST stacks, gadget state pages).\\
     PML4[0..255] map guest userspace via EPT-backed physical memory.};
\end{tikzpicture}
\caption{Address space layout for the \Seccomp{} backend.  The \KVM{}
backend additionally constructs guest-side PML4 page tables
that map the same physical memory through EPT/NPT.}
\label{fig:address-space}
\end{figure}

\subsection{Seccomp Backend Layout}

The \Seccomp{} backend's stub process has the following regions in its
virtual address space (Figure~\ref{fig:address-space}):

\begin{enumerate}
  \item \textbf{Guest userspace} (\texttt{0x0} to approximately
        \texttt{TASK\_SIZE}).  Mapped by the \UML{} kernel via batched
        \syscall{mmap}/\syscall{mprotect} calls into the stub process.

  \item \textbf{Stub pages} (\texttt{STUB\_START} to \texttt{STUB\_END}).
        Three pages: one page of stub executable code (the \texttt{SIGSYS}
        handler, the \syscall{futex} wait/wake loop, the syscall-batch
        executor) and two pages of \texttt{stub\_data} (shared state:
        registers, signal information, the futex word, and the
        syscall-batch buffer).

  \item \textbf{UML physical memory} (\texttt{uml\_physmem}).
        The contiguous region backing guest RAM, typically \syscall{mmap}'d
        from a \texttt{tmpfs} file descriptor or as anonymous memory.
        Guest page tables reference offsets within this region.

  \item \textbf{UML kernel text/data.}  The \UML{} binary's own code
        and data segments, loaded by the host's ELF loader.
\end{enumerate}

\subsection{KVM Backend: Guest Page Tables}

The \KVM{} backend constructs a full x86-64 PML4 page-table hierarchy
for the guest.  \UML{}'s physical memory is registered with \KVM{}
via \texttt{KVM\_SET\_USER\_MEMORY\_REGION}, which populates EPT/NPT
entries so that guest-physical addresses translate through host-virtual
addresses to host-physical pages without hypervisor intervention on
each access.

The guest PML4 contains a \emph{kernel half} at PML4 slot~448
(guest virtual address \texttt{0xffffe000\_00000000}) that maps the
trampoline page (LSTAR entry point), the IDT, exception handler stubs,
the GDT, per-vCPU TSS and IST stack pages, and per-vCPU gadget state
pages.  This kernel-half mapping is marked \texttt{US=0}
(supervisor-only) so guest userspace at CPL=3 cannot access it
directly; the CPU transitions to CPL=0 only transiently during the
\texttt{SYSCALL} instruction's microarchitectural handling before the
LSTAR trampoline code runs.


% ------------------------------------------------------------------
\section{Trap Mechanisms Compared}
\label{sec:foundations:trap-compare}

The performance of \UML{} is dominated by the cost of intercepting
guest system calls.  Three mechanisms have been used historically and
in the current implementation.  Figure~\ref{fig:trap-paths} illustrates
the data flow for each.

\subsection{ptrace: PTRACE\_SYSEMU}

The original \UML{} trap mechanism uses \syscall{ptrace} to intercept
every guest system call:

\begin{enumerate}
  \item Guest userspace executes the \texttt{syscall} instruction.
  \item The host kernel intercepts it (via \texttt{PTRACE\_SYSEMU})
        and stops the tracee with signal \texttt{SIGTRAP|0x80}.
  \item The \UML{} kernel (the tracer) wakes from \syscall{waitpid}.
  \item \texttt{PTRACE\_GETREGS} reads the tracee's register state.
  \item The \UML{} kernel dispatches the appropriate handler.
  \item \texttt{PTRACE\_SETREGS} writes the return value.
  \item \texttt{PTRACE\_SYSEMU} resumes the tracee.
\end{enumerate}

Each round-trip requires two context switches between the tracer and
the tracee plus multiple kernel transitions.  Measured latency:
\textbf{\SIrange{1}{5}{\micro\second}} per system call.

\subsection{seccomp: SIGSYS + futex}

The \Seccomp{} backend (contributed by Johannes Berg, circa 2022)
replaces the external-tracer model with an in-process signal:

\begin{enumerate}
  \item Guest userspace executes the \texttt{syscall} instruction.
  \item The host kernel's seccomp filter returns
        \texttt{SECCOMP\_RET\_TRAP}, delivering \texttt{SIGSYS} to
        the stub thread---\emph{without executing the system call}.
  \item The stub's \texttt{SIGSYS} handler saves register state into
        the shared \texttt{stub\_data} page.
  \item The stub issues \texttt{FUTEX\_WAKE} on the shared futex word,
        waking the \UML{} kernel thread.
  \item The \UML{} kernel reads registers from \texttt{stub\_data},
        dispatches the handler, and writes the result back.
  \item The \UML{} kernel issues \texttt{FUTEX\_WAKE}; the stub wakes
        from its \texttt{FUTEX\_WAIT} and resumes guest execution from
        the signal return path.
\end{enumerate}

Because no external tracer is involved, the context-switch overhead is
eliminated.  The signal delivery and futex round-trip are the dominant
costs.  Measured latency: \textbf{\SIrange{300}{500}{\nano\second}}
per system call---roughly an order of magnitude faster than \syscall{ptrace}.

\subsection{KVM: vmentry/vmexit + LSTAR Gadget}

The \KVM{} backend replaces the signal-based mechanism with hardware
VM entry/exit:

\begin{enumerate}
  \item Guest userspace executes the \texttt{syscall} instruction.
  \item The CPU loads \texttt{MSR\_LSTAR} into RIP (standard x86-64
        \texttt{SYSCALL} behavior).  Inside the guest VM, LSTAR points
        to the \UML{} trampoline page.
  \item The LSTAR trampoline executes \texttt{out \%al, \$0xf4}
        (an I/O instruction on port \texttt{0xf4}).
  \item The I/O instruction causes a \texttt{KVM\_EXIT\_IO} VM exit.
  \item The \UML{} kernel (running in VMX root mode via the
        \texttt{KVM\_RUN} return path) reads the exit reason, extracts
        registers from the shared \texttt{kvm\_run} page, and dispatches
        the system-call handler.
  \item The handler writes the return value into \texttt{kvm\_run},
        and the next \texttt{KVM\_RUN} call re-enters the guest.
        The guest resumes via \texttt{sysretq} from the trampoline.
\end{enumerate}

Measured latency for the full round-trip:
\textbf{\SIrange{100}{150}{\nano\second}} per system call.

\subsubsection{The LSTAR Gadget: Avoiding the Exit Entirely}

For a curated set of trivial system calls (\syscall{getpid},
\syscall{gettid}, \syscall{getppid}, \syscall{getuid},
\syscall{geteuid}, \syscall{getgid}, \syscall{getegid},
\syscall{getcpu}, \syscall{time}, and
\texttt{clock\_gettime(CLOCK\_MONOTONIC)}), the LSTAR trampoline
contains an expanded 343-byte \emph{gadget} that handles the system
call entirely within the guest, without any VM exit:

\begin{enumerate}
  \item \texttt{swapgs} switches GS base to a per-vCPU state page
        populated by the host on each context switch.
  \item A dispatch chain (\texttt{cmp}/\texttt{je} on the syscall
        number in \texttt{EAX}) routes to the appropriate handler.
  \item The handler reads the pre-computed return value from the
        per-vCPU state page via \texttt{\%gs:offset}.
  \item \texttt{swapgs} + \texttt{sysretq} returns directly to guest
        userspace.
\end{enumerate}

The entire sequence---including the \texttt{SYSCALL}/\texttt{SYSRET}
pair---completes in approximately \textbf{90~cycles} ($\approx$
\SI{30}{\nano\second} at \SI{3}{\giga\hertz}).  For workloads that
call \syscall{getpid} or \texttt{clock\_gettime} in tight loops, this
approaches bare-metal performance.

\begin{figure}[ht]
\centering
\begin{tikzpicture}[
  box/.style={draw, thick, rounded corners=3pt, minimum height=0.7cm,
              font=\footnotesize, text width=2.3cm, align=center},
  arrow/.style={-{Stealth[length=5pt]}, thick},
  phase/.style={font=\footnotesize\bfseries, color=UMLBlue},
  timing/.style={font=\scriptsize\itshape, color=UMLRed}
]
  % ptrace path
  \node[phase] at (-1.5, 5.5) {ptrace};
  \node[box, fill=KernelGold!15] (p1) at (0, 5) {Guest\\syscall};
  \node[box, fill=UMLRed!10]     (p2) at (3, 5) {Host kernel\\SYSEMU stop};
  \node[box, fill=UMLBlue!10]    (p3) at (6, 5) {waitpid\\GETREGS};
  \node[box, fill=UMLGreen!10]   (p4) at (9, 5) {UML\\handler};
  \node[box, fill=UMLBlue!10]    (p5) at (12, 5) {SETREGS\\SYSEMU};
  \draw[arrow] (p1) -- (p2);
  \draw[arrow] (p2) -- (p3);
  \draw[arrow] (p3) -- (p4);
  \draw[arrow] (p4) -- (p5);
  \draw[arrow, dashed] (p5) .. controls (14, 5) and (14, 5.8) ..
    (12.5, 5.8) -- (-0.5, 5.8) .. controls (-2, 5.8) and (-2, 5) .. (p1);
  \node[timing] at (6, 4.3) {$\sim$1--5 \textmu{}s per syscall};

  % seccomp path
  \node[phase] at (-1.5, 3.0) {seccomp};
  \node[box, fill=KernelGold!15] (s1) at (0, 2.5) {Guest\\syscall};
  \node[box, fill=UMLAmber!10]   (s2) at (3, 2.5) {SIGSYS\\handler};
  \node[box, fill=UMLTeal!10]    (s3) at (6, 2.5) {futex wake\\UML kernel};
  \node[box, fill=UMLGreen!10]   (s4) at (9, 2.5) {UML\\handler};
  \node[box, fill=UMLTeal!10]    (s5) at (12, 2.5) {futex wake\\stub};
  \draw[arrow] (s1) -- (s2);
  \draw[arrow] (s2) -- (s3);
  \draw[arrow] (s3) -- (s4);
  \draw[arrow] (s4) -- (s5);
  \draw[arrow, dashed] (s5) .. controls (14, 2.5) and (14, 3.3) ..
    (12.5, 3.3) -- (-0.5, 3.3) .. controls (-2, 3.3) and (-2, 2.5) .. (s1);
  \node[timing] at (6, 1.8) {$\sim$300--500 ns per syscall};

  % KVM path
  \node[phase] at (-1.5, 0.5) {KVM};
  \node[box, fill=KernelGold!15] (k1) at (0, 0) {Guest\\SYSCALL};
  \node[box, fill=UMLAmber!15]   (k2) at (3, 0) {LSTAR\\out \$0xf4};
  \node[box, fill=UMLRed!10]     (k3) at (6, 0) {KVM\_EXIT\\\_IO};
  \node[box, fill=UMLGreen!10]   (k4) at (9, 0) {UML\\handler};
  \node[box, fill=UMLBlue!10]    (k5) at (12, 0) {KVM\_RUN\\re-enter};
  \draw[arrow] (k1) -- (k2);
  \draw[arrow] (k2) -- (k3);
  \draw[arrow] (k3) -- (k4);
  \draw[arrow] (k4) -- (k5);
  \draw[arrow, dashed] (k5) .. controls (14, 0) and (14, 0.8) ..
    (12.5, 0.8) -- (-0.5, 0.8) .. controls (-2, 0.8) and (-2, 0) .. (k1);
  \node[timing] at (6, -0.7) {$\sim$100--150 ns per syscall
    (gadget: $\sim$90 cycles, no exit)};
\end{tikzpicture}
\caption{System call interception paths for the three \UML{} trap
mechanisms.  Dashed arrows show the return-to-guest path.}
\label{fig:trap-paths}
\end{figure}

\subsection{Comparison Table}

\begin{table}[ht]
\centering
\caption{Trap mechanism comparison.  Latencies are representative
measurements for a null system call (e.g., \syscall{getpid}) on
contemporary hardware (Intel Alder Lake / AMD Zen~4, 2024--2026).}
\label{tab:trap-comparison}
\small
\begin{tabularx}{\textwidth}{@{} l Y Y Y @{}}
\toprule
\textbf{Property} &
  \textbf{ptrace} &
  \textbf{seccomp} &
  \textbf{KVM (v2)} \\
\midrule
Trap mechanism &
  \texttt{PTRACE\_SYSEMU} $\rightarrow$ \texttt{SIGTRAP|0x80} &
  \texttt{SECCOMP\_RET\_TRAP} $\rightarrow$ \texttt{SIGSYS} &
  \texttt{SYSCALL} $\rightarrow$ LSTAR $\rightarrow$ \texttt{out} $\rightarrow$ \texttt{KVM\_EXIT\_IO} \\
\midrule
Syscall latency &
  \SIrange{1}{5}{\micro\second} &
  \SIrange{300}{500}{\nano\second} &
  \SIrange{100}{150}{\nano\second} \\
\midrule
Gadget fast-path &
  N/A &
  N/A &
  $\sim$90 cycles ($\sim$\SI{30}{\nano\second}), no VM exit \\
\midrule
Context switches &
  2 per syscall (tracer $\leftrightarrow$ tracee) &
  0 (in-process signal) &
  1 VM exit + 1 VM entry \\
\midrule
Host kernel requirement &
  \texttt{PTRACE\_SYSEMU} support &
  \texttt{SECCOMP\_RET\_TRAP} (kernel $\geq$ 3.5) &
  \texttt{/dev/kvm} + VT-x/AMD-V \\
\midrule
Guest memory model &
  \syscall{mmap} into stub process &
  \syscall{mmap} into stub process &
  \texttt{KVM\_SET\_USER\_MEMORY\_REGION} (EPT/NPT) \\
\midrule
SMP model &
  One host process per guest CPU &
  One host thread per guest CPU &
  One vCPU per host CPU (pool) \\
\midrule
Debugging support &
  Full ptrace attach &
  Limited (signal-based) &
  \texttt{KVM\_GUESTDBG\_*} (breakpoints, single-step) \\
\midrule
Kernel version in-tree &
  v2.4.x -- present (mainline) &
  v6.1 -- present (Berg series) &
  v2 in development (this work) \\
\bottomrule
\end{tabularx}
\end{table}


% ------------------------------------------------------------------
\section{Memory Management Under UML}
\label{sec:foundations:memory}

\subsection{Guest Physical Memory}

\UML{}'s ``physical memory'' is a contiguous host-virtual region of
size \texttt{physmem\_size} (set by the \texttt{mem=} boot parameter,
default \SI{64}{\mebi\byte}).  It is backed either by a file on
\texttt{tmpfs} (the traditional approach) or by an anonymous
\texttt{MAP\_PRIVATE} mapping.  The file-backed approach supports
\texttt{MAP\_SHARED} semantics, enabling host-side inspection of
guest memory; the anonymous approach avoids filesystem overhead.

Guest physical addresses are simple offsets into this region:
\begin{equation*}
  \text{host\_va}(\text{gpa}) = \texttt{uml\_physmem} + \text{gpa}
\end{equation*}

\subsection{Guest Page Tables}

\UML{} maintains full x86-64 four-level page tables (PML4 $\rightarrow$
PDP $\rightarrow$ PD $\rightarrow$ PT) in guest physical memory.  These
page tables serve different purposes depending on the backend:

\begin{description}[style=nextline]
  \item[Seccomp/ptrace backends.]
    Guest page tables exist only as a data structure internal to the
    \UML{} kernel.  They record the guest's virtual-to-physical
    mappings, but the \emph{actual} address translation is performed
    by the host kernel's page tables (the stub process's
    \texttt{mm\_struct}).  When the guest modifies a PTE, \UML{} must
    call \syscall{mmap}/\syscall{mprotect} on the stub process to
    make the host-side mapping match.

  \item[KVM backend.]
    Guest page tables are loaded into the vCPU's CR3 and walked by
    hardware during guest execution.  EPT/NPT provides the second
    level of translation (guest-physical to host-physical).  \UML{}
    still maintains the guest page tables as data structures, but
    they are now \emph{live}---the CPU reads them directly.
\end{description}

\subsection{The TLB Synchronization Problem}
\label{sec:foundations:tlb-sync}

When the \UML{} kernel modifies a guest page-table entry---due to a
page fault, an \syscall{mmap}, an \syscall{mprotect}, or a process
exit---the change must be propagated to the execution substrate:

\begin{description}[style=nextline]
  \item[Seccomp backend:]
    Each PTE change is translated into a host \syscall{mmap} or
    \syscall{mprotect} call, batched into the stub's syscall buffer
    (\texttt{stub\_data->syscall\_data}), and flushed on the next
    stub round-trip.  This is the \texttt{um\_tlb\_sync()} path in
    \umlpath{arch/um/kernel/tlb.c}.

  \item[KVM backend:]
    PTE changes are written directly into the guest page-table pages
    (which live in \UML{}'s physical memory region).  The \UML{} kernel
    then issues \texttt{INVLPG} or a full CR3 reload within the guest
    to invalidate stale TLB entries.  On SMP configurations,
    \texttt{tlb\_kick\_others()} sends an IPI to force all other vCPUs
    to flush their TLBs as well---this is the \KVM{}-specific backend
    op that the \Seccomp{} backend leaves NULL (the host kernel's
    mmu\_notifier mechanism handles cross-CPU invalidation
    transparently for the \Seccomp{} path).
\end{description}

The TLB sync path is one of the most performance-sensitive operations
in \UML{}.  The \Seccomp{} backend batches multiple PTE changes into
a single stub round-trip (amortizing the futex wake/wait overhead over
dozens of mappings), while the \KVM{} backend avoids host system calls
entirely for the common case, paying only the cost of a guest-side
\texttt{INVLPG} or CR3 write.


% ------------------------------------------------------------------
\section{Signal Handling}
\label{sec:foundations:signals}

Signals are the ``interrupts'' of the \UML{} kernel.  Where a
bare-metal kernel receives timer ticks via a local APIC interrupt and
I/O notifications via device IRQs, \UML{} receives:

\begin{description}[style=nextline]
  \item[\texttt{SIGALRM}:]
    Timer interrupts.  The \UML{} kernel programs a host
    \texttt{timer\_create}/\texttt{timer\_settime} timer (or
    \texttt{setitimer} on older configurations) that delivers
    \texttt{SIGALRM} at the frequency required by the guest's
    \texttt{HZ} setting.  This drives the scheduler tick, watchdog
    timers, and all time-dependent kernel subsystems.

  \item[\texttt{SIGIO}:]
    I/O readiness notifications.  \UML{}'s virtual device drivers
    register host file descriptors with the \UML{} IRQ subsystem,
    which uses \texttt{epoll} (or \texttt{poll}) and arranges for
    \texttt{SIGIO} delivery when any registered descriptor becomes
    ready.  This is the functional equivalent of a hardware IRQ
    line going active.

  \item[\texttt{SIGSEGV}:]
    Page faults.  When guest code accesses an address that is not
    mapped in the stub process's host-side page tables (for the
    \Seccomp{} backend) or that causes an EPT violation (for the
    \KVM{} backend), the resulting fault is delivered to the \UML{}
    kernel, which resolves it by consulting the guest page tables,
    allocating physical pages if necessary, and updating the mapping.

  \item[\texttt{SIGSYS}:]
    System call interception (\Seccomp{} backend only).  As described
    in Section~\ref{sec:foundations:trap-compare}, the seccomp filter's
    \texttt{SECCOMP\_RET\_TRAP} action delivers this signal to indicate
    that the guest issued a system call.
\end{description}

\subsection{Signal Delivery Across Backends}

The signal model differs substantially between backends:

\begin{itemize}
  \item \textbf{Seccomp backend:}
    \texttt{SIGALRM} and \texttt{SIGIO} are delivered asynchronously
    to the \UML{} kernel thread (the ``spawner'' process or the per-mm
    worker process).  The signal handler runs in the \UML{} kernel's
    context, interrupting whatever kernel code was executing.  Signal
    masking ensures that only one signal is handled at a time.

  \item \textbf{KVM backend:}
    While a vCPU is inside \texttt{KVM\_RUN}, host signals cannot be
    delivered to the thread---the thread is in kernel space executing
    the KVM entry/exit loop.  Instead, pending signals cause
    \texttt{KVM\_RUN} to return with \texttt{-EINTR} (or, with
    signal-window-based delivery, the signal is handled before the
    next \texttt{KVM\_RUN} call).  The \UML{} kernel's vCPU run loop
    checks for pending signals after each VM exit and dispatches
    them through the standard \UML{} interrupt-handling path.

    This means \texttt{SIGALRM} timer ticks are effectively ``held''
    during guest execution and processed at the next VM exit
    boundary.  For the LSTAR gadget fast path, which handles system
    calls without any VM exit, this creates a timer-delivery gap:
    a tight loop calling only gadget-handled system calls could
    theoretically starve the timer.  The gadget's \emph{budget}
    mechanism addresses this: each gadget-handled
    \texttt{clock\_gettime} call decrements a per-vCPU budget counter,
    and when the budget reaches zero, the gadget falls back to the
    full VM exit path, ensuring the host has an opportunity to deliver
    pending signals.
\end{itemize}

\subsection{Signal Masking and Reentrancy}

\UML{}'s signal handlers must be carefully written to avoid reentrancy
issues.  The \UML{} kernel uses \texttt{sigprocmask} to block signals
during critical sections (analogous to \texttt{local\_irq\_disable()}
on bare metal) and unblocks them at safe points (analogous to
\texttt{local\_irq\_enable()}).  The \texttt{block\_signals\_trace()} /
\texttt{unblock\_signals\_trace()} wrappers in
\umlpath{arch/um/os-Linux/signal.c} implement this protocol, and all
signal-handler dispatchers use them to ensure that a \texttt{SIGALRM}
arriving during \texttt{SIGIO} handling (or vice versa) is deferred
rather than creating a nested handler invocation that could corrupt
kernel state.

\bigskip

With these foundations established---the virtualization taxonomy, the
x86 privilege model, hardware virtualization extensions, the Linux
process model, address-space layout, trap mechanisms, memory
management, and signal handling---we are now equipped to examine the
architecture of \UML{} itself in Chapter~\ref{ch:architecture}.
